\section{Từ tiên đề đến lời giải thích cộng tính}
\label{sec:ch04-tu-tien-de-den-shap}

Trong lời giải thích mô hình, người chơi là các đặc trưng. Lundberg và Lee đặt
một \tn{lời giải thích cộng tính}{additive explanation} dưới dạng mô hình
đơn giản trên biểu diễn nhị phân $z'$:

\begin{equation}
\label{eq:ch04-mo-hinh-cong-tinh}
g(z')=\phi_0+\sum_{i=1}^{M}\phi_i z'_i.
\end{equation}

Ở đây $z'_i=1$ nghĩa là đặc trưng $i$ có mặt trong liên minh đang xét. Hệ số
$\phi_i$ là phần đóng góp được phân cho đặc trưng đó, còn $\phi_0$ là giá trị
nền. Dạng của phương trình~\ref{eq:ch04-mo-hinh-cong-tinh} giống mô hình
tuyến tính, nhưng các hệ số không phải trọng số học trực tiếp từ dữ liệu.
Chúng là kết quả của việc đánh giá các liên minh rồi áp công thức Shapley.

Trong lớp mô hình này, bài báo SHAP dùng ba tính chất gần với các tiên đề trò
chơi. \tn{Độ chính xác cục bộ}{local accuracy} yêu cầu tổng giá trị nền và các đóng góp khớp
đúng đầu ra tại $x$. \tn{Tính vắng mặt}{missingness} yêu cầu một đặc trưng không có mặt trong
biểu diễn không nhận đóng góp. \tn{Tính nhất quán}{consistency} yêu cầu đóng góp của một
đặc trưng không giảm nếu ảnh hưởng biên của nó không giảm trong mọi liên minh.
Với một cách ánh xạ biểu diễn đã cố định, ba tính chất ấy xác định duy nhất
các \tn{giá trị Shapley}{Shapley value}~\autocite{p10shap}.

Do đó, SHAP không nói rằng mọi lời giải thích cộng tính đều trung thực. Nó
nói một điều hẹp hơn và mạnh hơn: nếu đã chọn mô hình giải thích cộng tính cùng
quy tắc ánh xạ, thì ba tính chất trên chỉ còn một cách phân chia đóng góp. Câu
hỏi còn lại là giá trị của một liên minh phải được định nghĩa từ $f$ và dữ liệu
ra sao.

\index{lời giải thích cộng tính}%
